\documentclass{article}
\usepackage{graphicx} % Required for inserting images
\usepackage{amsmath}
\usepackage{amsfonts}
\usepackage{listings}
\usepackage{booktabs}
\usepackage{float}
\usepackage{hyperref} % should always go last

\renewcommand{\labelenumi}{\alph{enumi})}

\lstset{
  basicstyle=\ttfamily\footnotesize,
  frame=single,
  captionpos=t,
  breaklines=true,
}

\newcommand{\includesource}[1]{%
  \lstinputlisting[caption={\texttt{#1}}]{#1}%
}

\graphicspath{{1/}}   % the four PDFs live in the folder "1"

\title{Performance Lab 7 - Group 8}
\author{Hans Bouwman, Joel de Haan}

\begin{document}

\maketitle

\section{Data Structure Optimisation}

For \texttt{Small}, padding causes its size to be 72 bytes, moving \texttt{flg2} before the array reduces its size to 64 bytes, which combined with an alignas makes sure it fits snugly in a cacheline. Packing would only destroy the cache alignment and would only reduce the size to 56, which is probably not worth it.

For \texttt{Big}, reversing the order would give the same size and packing would only save 4 bytes relative to $2^14$. Similarly, the data is very large so padding could not help the structure fit into a cache line. The only optimisation is to split the rank from the array: this way we can sort on the ranks (introduced with a new index fields) and implicitly access elements of the array.

Thus we propose the following optimised data structure:

\begin{lstlisting}
struct alignas(64) Small
{
    char flg1;
    char flg2;
    double x[7];
};

struct BigMeta
{
    int rank;
    int arrayIndex;
};

struct BigArray
{
    double x[elementCount];
};
\end{lstlisting}


As seen in table \ref{tab:perf-small}, there is a 21\% reduction in loads because of the reduced size and a 13\% reduction in load misses due to better cache alignment, which lines up very well with the 14\% reduction in total time taken (cache misses are of course a major source of slowdown).

In table \ref{tab:perf-big}, we see an enormous reduction in loads, as is expected from splitting, and similarly a 20 fold reduction in cache miss rate, likely due to the fact that multiple vector elements now fit in one cache line.



\begin{table}[htbp]
\centering
\begin{tabular}{lrrr}
\toprule
Counter & Original & Optimized & Original/Optimized \\
\midrule
Time per sort (ms) & 0.187 & 0.163 & 1.14 \\
Cycles & $6.79\cdot10^{5}$ & $5.98\cdot10^{5}$ & 1.13 \\
Instructions & $1.73\cdot10^{6}$ & $1.62\cdot10^{6}$ & 1.07 \\
Branch misses & 843 & 841 & 1 \\
L1d loads & $8.30\cdot10^{5}$ & $6.84\cdot10^{5}$ & 1.21 \\
L1d load misses & $1.18\cdot10^{5}$ & $1.04\cdot10^{5}$ & 1.13 \\
L1d miss rate (\%) & 14.2 & 15.2 & 0.933 \\
\bottomrule
\end{tabular}
\caption{Small: perf counters per sort, original layout versus reordered fields and \texttt{alignas(64)} ($2^{14}$ elements). Means over 3 runs with counters enabled only during the sorts (perf \texttt{--control}).}
\label{tab:perf-small}
\end{table}

\begin{table}[htbp]
\centering
\begin{tabular}{lrrr}
\toprule
Counter & Original & Optimized & Original/Optimized \\
\midrule
Time per sort (ms) & 153 & 0.048 & $3.18\cdot10^{3}$ \\
Cycles & $5.88\cdot10^{8}$ & $1.42\cdot10^{5}$ & $4.13\cdot10^{3}$ \\
Instructions & $2.34\cdot10^{8}$ & $5.30\cdot10^{5}$ & 443 \\
Branch misses & $4.00\cdot10^{4}$ & 404 & 99.1 \\
L1d loads & $2.30\cdot10^{8}$ & $1.57\cdot10^{5}$ & $1.47\cdot10^{3}$ \\
L1d load misses & $7.59\cdot10^{7}$ & $2.64\cdot10^{3}$ & $2.87\cdot10^{4}$ \\
L1d miss rate (\%) & 33 & 1.69 & 19.5 \\
\bottomrule
\end{tabular}
\caption{Big: perf counters per sort, original layout versus split into \texttt{BigMeta} and \texttt{BigArray} ($2^{13}$ elements of $2^{13}$ doubles). Means over 3 runs with counters enabled only during the sorts (perf \texttt{--control}).}
\label{tab:perf-big}
\end{table}


\newpage

\section{MPI Benchmark}

The bandwith is given by 
$$\frac{2N R}{T} \text{Bytes/s}$$
where $N$ is the message size (in Bytes) and $R$ the number of repetitions.

The experiment was performed incrementing N in powers of 2 from 128 to 4096 and 1000 repetitions.
For the single node this give the following values:
7.746e+07
1.26484e+08
2.34318e+08
4.17738e+08
6.58821e+08
1.00676e+09

which have an average of 4.20264e+08.


For the 2 nodes the values are:
5.88728e+08
9.95685e+08
1.66794e+09
2.08036e+09
3.12324e+09
4.3928e+09

where it can be seen that the volatility is higher and has an average of 2.14146e+09.




\end{document}
